\subsection{Du catalogue des réalisations à l'inventaire de 471 masses}
L'extension universelle suit cinq opérations :

\[
\begin{aligned}
 \text{identité physique}
 &\longrightarrow
 \text{projecteur de représentation}
 \longrightarrow
 \text{registre généalogique élémentaire ou composé}\\
 &\longrightarrow
 \text{opérateur/pôle}
 \longrightarrow
 \text{carte de comparaison}.
\end{aligned}
\]
Pour les états élémentaires, la deuxième application agit sur une fibre de l'atlas. Pour les mésons et les baryons, elle agit sur les 432 et 1728 expressions composées. Pour les résonances, la quatrième application donne un pôle complexe. Pour les secteurs topologiques, elle donne d'abord un gap spectral et une représentation de tresse. Ce typage permet de couvrir tout l'inventaire sans déclarer que chaque nom externe correspond à une cellule élémentaire.

\Needspace{7\baselineskip}
\subsection{Des 384 largeurs à l'opérateur de survie}
Chaque largeur exige

\[
 X\longmapsto Q_XUQ_X
 \longmapsto \rho(Q_XUQ_X)
 \longmapsto\Gamma_X
\]
ou, de manière équivalente, le pôle de \(H_{\rm eff}\). Le lecteur doit incorporer canaux de désintégration, multiplicité, seuil et espace des phases. L'égalité

\[
 \Gamma_{\bar X}=\Gamma_X
\]
ne se déduit que si l'opérateur de désintégration est covariant par conjugaison.

\Needspace{7\baselineskip}
\section{Masse, température et statistique}
\Needspace{7\baselineskip}
\subsection{Transducteur thermique entre caractère de masse, température et occupation statistique}
La constante de Boltzmann est le transducteur

\[
 k_B:\mathcal L_\Theta\longrightarrow\mathcal L_E.
\]
L'horloge CT--108 fixe l'identité

\[
 k_B\Theta_{\rm clk}\log3
 =\hbar_{\rm int}\frac{2\pi}{108t_0}.
\]
Cette équation relie une unité d'information tritique, une fréquence et une énergie. Elle n'introduit pas la coordonnée SI de \(k_B\) dans la loi ; elle détermine le produit typé une fois la carte thermique choisie.

\Needspace{7\baselineskip}
\subsection{Projection fermionique de l'atlas et distribution de Fermi–Dirac}
Pour les niveaux

\[
 \epsilon=(0,3,6,9),\qquad
 g=(1,6,12,8),
\]
et contraintes

\[
 (N,E)=(12,66),
\]
la complétion extérieure produit

\[
 2\,104\,494
\]
microétats fermioniques compatibles. Les multiplicateurs effectifs sont

\[
 (\beta,\mu)
 =(0.1530701135,\;4.5028651976).
\]
Cette réalisation illustre comment masse, énergie et exclusion émergent des parcours fermioniques sans que la statistique soit un ajustement supplémentaire.

\Needspace{7\baselineskip}
\subsection{Projection bosonique de l'atlas et distribution de Bose–Einstein}
Sur les mêmes niveaux et sous les mêmes contraintes, la complétion symétrique produit

\[
 259\,436\,430
\]
microétats, avec

\[
 (\beta,\mu)
 =(0.0545672796,\;-15.9256978191).
\]
La différence entre Fermi--Dirac et Bose--Einstein est donc une différence dans la règle de composition des parcours, et non dans la valeur d'une masse individuelle.

\Needspace{7\baselineskip}
\subsection{Critère de condensation des sections bosoniques et occupation macroscopique}
Un condensat bosonique apparaît lorsqu'une composante symétrique acquiert une occupation macroscopique dans la valeur propre minimale de l'opérateur. Une mer fermionique, en revanche, remplit un ensemble de valeurs propres distinctes jusqu'à une frontière. HMT décrit les deux au moyen du même spectre de masse/énergie et de deux foncteurs de Fock distincts. Cette séparation empêche d'identifier « beaucoup de particules sur un parcours » à une particule plus massive.

\Needspace{7\baselineskip}
\section{Propagation différentielle des incertitudes et stabilité structurelle du caractère de masse}
\Needspace{7\baselineskip}
\subsection{Invariance du graphe génératif sous la confrontation métrologique ultérieure}
Le résultat interne est figé dans l'ordre

\[
 \mathcal R_{324}
 \to P_X
 \to\mathcal L_X
 \to\sigma_X
 \to\eta_X
 \to\widehat M_X
 \to\mu_X.
\]
Ce n'est qu'ensuite que la carte externe est ouverte. Une permutation arbitraire des 471 masses ou des 384 largeurs doit laisser invariants tous les objets précédents.

\Needspace{7\baselineskip}
\subsection{Décomposition de l'incertitude structurelle en contributions expérimentale, de carte, de parcours, de tour et de couplage}
L'incertitude totale se décompose comme

\[
 \Sigma_{\rm tot}
 =\Sigma_{\rm ext}
 +\Sigma_{\rm carta}
 +\Sigma_{\rm route}
 +\Sigma_{\rm tower}
 +\Sigma_{\rm coupling}.
\]
Ici :

\begin{itemize}
\item \(\Sigma_{\rm ext}\) est la covariance expérimentale ;
\item \(\Sigma_{\rm carta}\) quantifie schéma, échelle et conversion des unités ;
\item \(\Sigma_{\rm route}\) compare tous les parcours de l'orbite admissible ;
\item \(\Sigma_{\rm tower}\) utilise la borne de reste de la série harmonique ;
\item \(\Sigma_{\rm coupling}\) mesure la variation entre les entrelaceurs autorisés.
\end{itemize}
Elles ne s'additionnent pas nécessairement comme des nombres indépendants ; l'équation exprime une décomposition des sources qui doit être convertie en covariance conjointe.

\Needspace{7\baselineskip}
\subsection{Dérivée de l'observable de masse par rapport à la signature et aux lecteurs}
Pour

\[
 \log m_X=
 \log m_\star+
 \log\chi_0(\sigma_X)
 +\kappa_X\log R_{\rm act}
 +\sum_h\eta_{X,h}s_h,
\]
la variation est

\[
 d\log m_X
 =d\log m_\star
 +\langle\Theta_0,d\sigma_X\rangle
 +\langle\sigma_X,d\Theta_0\rangle
 +\log R_{\rm act}\,d\kappa_X
 +\kappa_X\,d\log R_{\rm act}
 +\sum_h(s_h\,d\eta_{X,h}+\eta_{X,h}\,ds_h).
\]
Cette formule permet d'attribuer chaque décimale au parcours, à une constante, au lecteur, à la tour ou à la carte. Un résultat est robuste lorsque les perturbations préservant les prémisses maintiennent sa classe spectrale et son incertitude calculée.

\Needspace{7\baselineskip}
\subsection{Essais de dépendance structurelle par suppression contrôlée du feuillet, de l'orientation, de la retenue, de la mémoire, de l'enroulement et de la frontière}
Un test de dépendance élimine séparément feuillet, orientation, retenue, mémoire, enroulement ou frontière, puis recalcule l'objet. Le résultat attendu n'est pas que tout change, mais que changent exactement les coordonnées dépendant de la composante éliminée. Si la suppression de la mémoire laisse invariant un coefficient défini par le retour, la construction de ce coefficient est réfutée.

\Needspace{7\baselineskip}
\subsection{Obstructions à l'identification des parcours, des caractères et des grandeurs métrologiques}
La loi est réfutée dans une application concrète si l'un des faits suivants se produit :

\begin{enumerate}
\item une masse externe modifie le parcours ou le projecteur choisi ;
\item la signature présentée n'est pas régénérée à partir de l'enregistrement généalogique ;
\item le raffinement ne converge pas ou dépend d'une troncature choisie selon le résidu ;
\item deux parcours équivalents reçoivent des masses distinctes sans brisure de symétrie ;
\item une antiparticule acquiert une masse distincte sans violation déclarée de la conjugaison ;
\item un triplet de couleur se scinde sans brisure de couleur ;
\item un scalaire est présenté alors que la variance spectrale est positive ;
\item une largeur est obtenue à partir du caractère de masse sans opérateur temporel ;
\item une carte de quark omet le schéma ou l'échelle ;
\item une instabilité est interprétée comme un tachyon persistant.
\end{enumerate}
\Needspace{7\baselineskip}
